% ============================================
% Johns Hopkins Letter Template
% ============================================

\documentclass[11pt,letterpaper]{letter}

\usepackage{graphicx}
\usepackage{eso-pic}
\usepackage{geometry}
\usepackage{microtype}
\usepackage[utf8]{inputenc}
\usepackage[hidelinks]{hyperref}

\geometry{left=25mm,right=25mm,top=25mm,bottom=25mm}

\newcommand{\PutJHULogoFG}[1]{%
  \AddToShipoutPictureFG*{%
    \AtPageUpperLeft{%
      \hspace{10mm}%
      \raisebox{-38mm}{%
        \includegraphics[width=6cm]{#1}%
      }%
    }%
  }%
}

\signature{Decory Edwards}
\address{
Department of Economics \\
Johns Hopkins University \\
Baltimore, MD 21218 \\
\texttt{dedwar65@jh.edu}
}

\begin{document}
\PutJHULogoFG{Images/jhulogo.png}

\begin{letter}{
Hiring Committee \\
Research Department \\
Federal Reserve Bank of Richmond
}

\opening{Dear Members of the Hiring Committee,}

I am writing to apply for the Post-Doc Economist position in the Research Department at the Federal Reserve Bank of Richmond. I am a Ph.D.\ candidate in Economics at Johns Hopkins University (advisors: Christopher D.\ Carroll, Nicholas W.\ Papageorge, and Stelios Fourakis), and I expect to complete my degree in May 2026. My primary research fields are macroeconomics, monetary economics, and financial economics, with an emphasis on heterogeneous agent modeling and quantitative policy analysis.

My research agenda focuses on the ability of macroeconomic models to generate empirically realistic distributions of wealth and income over the life cycle. A key driver of that work is modeling heterogeneity in the rate of return on assets, and understanding how return dispersion contributes to portfolio accumulation across households.

In my job market paper, I calibrate a life cycle heterogeneous agent model to match wealth moments from the Survey of Consumer Finances by allowing households to differ in their portfolio returns. The model helps explain observed wealth concentration and return dispersion. In a related research project using the Health and Retirement Study, I examine how trust influences both income growth and asset returns. I find a hump shaped relationship for trust and returns and characterize the distribution of the persistent component of returns to net worth.

Beyond publication-oriented research, my work is inherently policy relevant. By studying how heterogeneity shapes household responses to shocks, my research speaks directly to questions concerning monetary transmission, financial market participation, and distributional consequences of policy. I have extensive experience using Python and Stata to manage large survey datasets, conduct structural estimation, and simulate counterfactual policy environments. I am comfortable presenting technical results to diverse audiences and translating complex quantitative findings into clear economic narratives.

I am particularly excited about the opportunity to contribute to the Richmond Fed’s research mission while engaging directly with policy questions relevant to the Fifth District and the broader Federal Reserve System. I would welcome the opportunity to conduct frontier research aimed at publication in leading journals, present my work at System meetings and academic conferences, and collaborate with other economists in the department. I am also enthusiastic about preparing policy briefings for Bank leadership and contributing to publications that inform internal and public economic perspectives. My training in macroeconomic modeling, econometrics, and financial economics provides a strong foundation for advising on issues related to monetary policy, banking and financial markets, and economic stability.

I am enthusiastic about working in person in Richmond and contributing to a collaborative research environment. I value institutions that combine academic rigor with public service, and I am motivated by the Federal Reserve’s mission to strengthen and protect the economy and the communities it serves.

I believe I would be a strong fit for this role because:
\begin{enumerate}
    \item I conduct rigorous, publication-quality research in macroeconomics and financial economics.
    \item My work is directly relevant to monetary policy and financial market analysis.
    \item I am committed to clear communication of complex economic issues to policymakers and leadership.
\end{enumerate}

Thank you very much for your consideration. I would welcome the opportunity to contribute to the research and policy mission of the Federal Reserve Bank of Richmond.

\closing{Sincerely,}

\end{letter}
\end{document}